\begin{frame}[fragile]{Eigen Commando's}
    Je kunt je \textbf{eigen commando's} aanmaken om veel typwerk te voorkomen. Je doet dit als volgt:
    \begin{minted}{tex}
        \newcommand{\naam}{wat je wil dat je commando doet}
    \end{minted}
    \pause 
    Je zet \mintinline{tex}|\newcommand| meestal \emph{voor} je \mintinline{tex}{\begin{document}}:
    \begin{minted}[linenos, highlightlines={2,4}]{tex}
        \documentclass[a4paper]{article}
        \newcommand{\R}{\mathbb{R}}
        \begin{document} 
        We kunnen nu ons nieuwe commando aanroepen: $x \in \R$. 
        Dit is sneller dan $x \in \mathbb{R}$
        \end{document}
    \end{minted}
    \newcommand{\mbbR}{\mathbb{R}}
    \pause
    \begin{tcolorbox}[boxrule=0.5pt, colback=white]
        We kunnen nu ons nieuwe commando aanroepen: $x \in \mbbR$. 
        Dit is sneller dan $x \in \mathbb{R}$
    \end{tcolorbox}
\end{frame}


%% Dit vind ie stom voor een of andere reden, sad

% \begin{frame}[fragile=singleslide]{Eigen Commando's}
%     Je kunt je eigen commando's aanmaken om veel typwerk te voorkomen.
%     \begin{columns}
%         \begin{column}{0.4\textwidth}
%             \begin{minted}[linenos, highlightlines={2,4}]{tex}
%                 \documentclass[a4paper]{article}
%                 \newcommand{\mbbR}{\mathbb{R}}
%                 \begin{document} 
%                 We kunnen nu ons nieuwe commando aanroepen: $x \in \mbbR$. 
%                 Dit is sneller dan $x \in \mathbb{R}$
%                 \end{document}
%             \end{minted}
%         \end{column} 
%         \newcommand{\mbbR}{\mathbb{R}}
%         \begin{column}{0.4\textwidth}
%             \begin{tcolorbox}[boxrule=0.5pt, colback=white]
%                 We kunnen nu ons nieuwe commando aanroepen: $x \in \mbbR$. 
%                 Dit is sneller dan $x \in \mathbb{R}$
%             \end{tcolorbox}
%         \end{column}
%     \end{columns}
% \end{frame}